% 张三慧第4版V1第8章：狭义相对论基础；独立计算候选解。
% 源教材物理页283--285，标准c=299792458m/s。
\section*{第8章\quad 狭义相对论基础：独立解析}

\paragraph{习题 8.1}
记 \(\beta=u/c\)，
\(\gamma=(1-\beta^2)^{-1/2}\)。
杆在静止系 \(S\) 的端点坐标差为
\(\Delta x=l\cos\theta,\quad\Delta y=l\sin\theta\)。
必须在 \(S'\) 中选择\emph{同一 \(t'\)}
的两个端点位置，才可测量运动长度；
沿相对运动方向收缩，垂直方向不变：
\[
 \Delta x'=\frac{l\cos\theta}{\gamma},\qquad
 \Delta y'=l\sin\theta.
\]
因此
\[
 \boxed{l'=l\sqrt{1-\beta^2\cos^2\theta}},\qquad
 \boxed{\tan\theta'=\gamma\tan\theta}
\]
（角度象限按端点顺序确定）。
只把总长度整体乘上 \(1/\gamma\)
会错误压缩横向分量。

\paragraph{习题 8.2}
把立方体一条边取在 \(x\) 方向，
静止时尺寸 \(a,a,a\)。
在 \(S'\) 同时测量，
运动方向边缩短至 \(a/\gamma\)，
其他两条边仍为 \(a\)，所以
\[
 \boxed{V'=\frac{a^3}{\gamma}
 =a^3\sqrt{1-u^2/c^2}}.
\]
它讨论同时空间截面的体积，
不是四维时空体积。

\paragraph{习题 8.3}
两枪固定在 \(S\) 系，间距为 \(L\)；
发射事件在 \(S\) 中同时
\(\Delta t=0,\ \Delta x=L\)。
在 \(S'\) 中它们的坐标差由
洛伦兹变换给出
\[
 \Delta x'=\gamma(\Delta x-u\Delta t)=\gamma L,\qquad
 \Delta t'=\gamma\left(\Delta t-
 \frac{u\Delta x}{c^2}\right)
 =-\frac{\gamma uL}{c^2}.
\]
两标记是刻在\emph{静止于 \(S'\)} 的尺上，
其两个刻度位置相距
\[
 \boxed{\Delta \ell'=\gamma L}.
\]
相差的是两个\emph{不同时刻}的击中事件，
而非 \(S'\) 中对运动的两枪作同时测长。
\(S'\) 观察者不会错误地以「同时测长」
要求两刻痕相距 \(L/\gamma\)。
这正说明相对论同时性的必要性。

\paragraph{习题 8.4}
在 \(S\) 中两事件同地点，
\(\Delta x=0,\Delta t=2\,\mathrm s\)；
\(S'\) 中
\(\Delta t'=\gamma\Delta t=3\,\mathrm s\)，
故 \(\gamma=3/2\)，
\(\beta=\sqrt{1-\gamma^{-2}}=\sqrt5/3\)。
空间间隔
\[
 \Delta x'=\gamma(\Delta x-u\Delta t)
 =-\gamma u(2\,\mathrm s)
 =-\sqrt5\,c(1\,\mathrm s),
\]
因此
\[
 \boxed{|\Delta x'|=\sqrt5\,c\,\mathrm s
 \approx6.70\times10^8\,\mathrm m}.
\]
符号依 \(S'\) 朝哪方向运动而定。

\paragraph{习题 8.5}
取两事件在 \(S\) 中同一时刻，
\(\Delta t=0,\Delta x=1\,\mathrm m\)。
若 \(S'\) 测得
\(|\Delta x'|=2\,\mathrm m\)，
由 \(|\Delta x'|=\gamma|\Delta x|\) 得
\(\gamma=2\)，
\(\beta=\sqrt3/2\)。
时间间隔
\[
 \Delta t'=-\gamma\frac{u\Delta x}{c^2};
\quad
 \boxed{|\Delta t'|=\frac{\sqrt3}{c}
 (1\,\mathrm m)\approx5.78\times10^{-9}\,\mathrm s}.
\]
该题给的 \(\Delta x'\)
是事件坐标差，不是收缩杆长度。

\paragraph{习题 8.6}
选地球系 \(S\)，飞船从地球沿 \(+x\)
以 \(\beta=4/5\)、\(\gamma=5/3\)
匀速远离；发信瞬间在地球系距地球
\(d_0\)，并取此刻 \(t=0\)。
信号到地球的时间为
\(t_1=d_0/c\)；
反射后返回，追赶飞船。
接收时刻 \(t_2\) 满足
\[
 c(t_2-t_1)=d_0+ut_2
 \Longrightarrow t_2=\frac{2d_0}{c-u}.
\]
发、收事件均发生于飞船世界线，
所以飞船钟的固有间隔
\(\Delta\tau=t_2/\gamma=60\,\mathrm s\)，
故 \(t_2=100\,\mathrm s\)。
\textbf{第（1）问要特别注意参考系：}
在飞船系 \(S'\) 中，
发送和收到信号发生在同一位置，
均为飞船原点，
固有时间间隔为 \(60\,\mathrm s\)；
因为光信号以速率 \(c\) 传播，
从飞船到地球反射位置、再返回飞船
的两段路程在这个系里都为
\(c(60/2)\,\mathrm s\)。
所以在\emph{地球反射时刻，飞船系}
测得的距离是
\[
 \boxed{d_{\rm reflection}^{(S')}
 =30c\,\mathrm s=9.00\times10^9\,\mathrm m}.
\]
作为中间校验，地球系在
\emph{发送瞬间}的距离为
\[
 d_{\rm send}^{(S)}=\frac12(c-u)t_2
 =10c\,\mathrm s=3.00\times10^9\,\mathrm m,
\]
但\textbf{这不是原题第（1）问所要求的距离}。
收信瞬间在地球系测得
\[
 \boxed{d_{\rm receive}^{(S)}
 =d_{\rm send}^{(S)}+ut_2
 =90c\,\mathrm s
 =2.70\times10^{10}\,\mathrm m}.
\]
配套辅导书PDF物理页127明确分别给出
\(9\times10^9\,\mathrm m\) 和
\(2.7\times10^{10}\,\mathrm m\)；
首次候选误把第（1）问转写为
发射时地球系距离，现已恢复题意。
不能将60秒直接当作地球系往返时间。

\paragraph{习题 8.7}
约定地球 \(S\) 坐标系、飞船 \(S'\)
沿 \(+x\) 匀速运动，
\(\beta=0.8,\gamma=5/3\)。
事件逆变换为
\[
 \boxed{t=\gamma\left(t'+\frac{ux'}{c^2}\right)},
 \quad
 \boxed{x=\gamma(x'+ut')},\quad
 y=y',\ z=z'.
\]
代入题给数值，按 \(c\simeq3.00\times10^8\,\mathrm{m/s}\)：
\[
 \boxed{t=-2.00\times10^8\,\mathrm s},\quad
 \boxed{x=6.00\times10^{16}\,\mathrm m},\quad
 \boxed{y=1.20\times10^{17}\,\mathrm m,\ z=0}.
\]
爆发位置相对飞船航向（地球系 \(+x\)）
的几何角：
\[
 \boxed{\alpha=\arctan(y/x)
 =\arctan 2\approx63.4^\circ}.
\]
事件的 \(t<0\) 表示爆发早于
飞船经过地球那一参考时刻，
不意味着信号早于爆发被接收。

\paragraph{习题 8.8}
在飞船系 \(S'\) 中，爆发位置距原点
\[
 r'=\sqrt{x'^2+y'^2}
 =\sqrt{1.8^2+1.2^2}\times10^{17}
 \approx2.1633\times10^{17}\,\mathrm m.
\]
来自超新星的光在
\(t'_{\rm arr}=t'_{\rm event}+r'/c\)
到达飞船，故
\[
 \boxed{t'_{\rm arr}\approx
 -6.0\times10^8+7.21\times10^8
 =1.21\times10^8\,\mathrm s}.
\]
宇航员立即将报告从其当前飞船位置
发向地球。
发报告的事件在飞船系原点
\(x'_{\rm arr}=0\)，
逆变换到地球系：
\(t_{\rm send}=\gamma t'_{\rm arr}\)，
\(x_{\rm send}=u t_{\rm send}>0\)。
地球接收报告还需时间
\(x_{\rm send}/c\)，所以
\[
 \boxed{t_{\rm report}=
 \gamma t'_{\rm arr}(1+\beta)
 \approx3.65\times10^8\,\mathrm s}.
\]
地球直接见到超新星的时间为
\[
 t_{\rm direct}=t_{\rm event}
 +\frac{\sqrt{x^2+y^2}}c,
\quad
 \boxed{t_{\rm direct}\approx2.47\times10^8\,\mathrm s}.
\]
为与源书的有效位数保持一致，以上用
\(c\simeq3.00\times10^8\,\mathrm{m/s}\)；
严格用定义值时报告到达时刻约
\(3.65\times10^8\,\mathrm s\)。
该题需要区别\emph{事件坐标时刻}、
\emph{光的到达时刻}和\emph{回报接收时刻}。

\paragraph{习题 8.9}
令飞船地球系速度 \(v_s=+0.60c\)，
彗星为 \(v_c=-0.80c\)。
飞船系的彗星速度按相对论速度合成：
\[
 v_c'=\frac{v_c-v_s}{1-v_cv_s/c^2}
 =\frac{-1.40c}{1+0.48}
 =-\frac{35}{37}c.
\]
因此
\[
 \boxed{|v_c'|=\frac{35}{37}c\approx0.946c}.
\]
碰撞事件位于飞船世界线上，
地球系现在至碰撞时间为5s；
飞船固有时间缩短：
\[
 \boxed{\Delta\tau=5\sqrt{1-0.60^2}
 =4.0\,\mathrm s}.
\]
不能把非相对论的接近速率 \(1.4c\)
当作飞船本系可测速度。

\paragraph{习题 8.10}
在 \(S'\) 中光的速度分量
\(v_x'=c\cos\theta'\)，
\(v_y'=c\sin\theta'\)，
变换到 \(S\)：
\[
 v_x=\frac{c\cos\theta'+u}
 {1+(u/c)\cos\theta'},\qquad
 v_y=\frac{c\sin\theta'}
 {\gamma[1+(u/c)\cos\theta']}.
\]
故传播方向角可由
\[
 \boxed{\tan\theta=
 \frac{\sin\theta'}{\gamma(\cos\theta'+u/c)}}
\]
确定，象限以 \(v_x,v_y\) 同时判定。
令 \(\beta=u/c\)，直接计算：
\[
 \frac{v_x^2+v_y^2}{c^2}
 =\frac{(\cos\theta'+\beta)^2
 +(1-\beta^2)\sin^2\theta'}
 {(1+\beta\cos\theta')^2}=1.
\]
所以 \(\boxed{|\boldsymbol v|=c}\)。

\paragraph{习题 8.11}
一维速度变换
\(v'=(v-u)/(1-uv/c^2)\)，
时间微分
\(dt'=\gamma(1-uv/c^2)\,dt\)。
先对 \(v\) 求导：
\[
 \frac{dv'}{dv}=
 \frac{1-u^2/c^2}{(1-uv/c^2)^2}.
\]
由 \(a'=dv'/dt'\)：
\[
 \boxed{a'=
 \frac{(1-u^2/c^2)^{3/2}}
 {(1-uv/c^2)^3}\,a}.
\]
低速极限
\(|u|/c,|v|/c\ll1\) 时
\(a'\to a\)，与伽利略变换下
惯性系加速度不变的结论一致。

\paragraph{习题 8.12}
沿力方向的相对论动量满足
\(dp/dt=F\)，
从静止出发给
\(p(t)=Ft\)。
由
\(p=\gamma m_0v\)：
\[
 \boxed{v(t)=
 \frac{Ft/m_0}{\sqrt{1+(Ft/(m_0c))^2}}}.
\]
位移从0积分：
\[
 \boxed{x(t)=\frac{m_0c^2}{F}
 \left[\sqrt{1+\left(\frac{Ft}{m_0c}\right)^2}-1\right]}.
\]
短时 \(Ft/(m_0c)\ll1\)，
\[
 \boxed{v\simeq Ft/m_0,\quad
 x\simeq Ft^2/(2m_0)}.
\]
长时 \(Ft/(m_0c)\gg1\)，
\[
 \boxed{v\to c^-,
 \quad x\simeq ct-\frac{m_0c^2}{F}
 +O(t^{-1})}.
\]
相对论动量可以无限增长，
但速率严格小于光速。


\paragraph{习题 8.13}
给定同一静质量 \(m_0\) 与速率 \(v\)，
非相对论动量 \(p_N=m_0v\)，
相对论动量 \(p_R=\gamma m_0v\)。
若 \(p_R=2p_N\)，则 \(\gamma=2\)：
\[
 \boxed{v=c\sqrt{1-\gamma^{-2}}
 =\frac{\sqrt3}{2}c\approx0.866c}.
\]
动能分别为
\(K_N=\tfrac12m_0v^2\)，
\(K_R=(\gamma-1)m_0c^2\)。
若 \(K_R=2K_N\)，则
\[
 \gamma-1=\beta^2
 =1-\frac1{\gamma^2}
 \Longrightarrow
 \gamma^2=\gamma+1.
\]
取正根
\(\gamma=(1+\sqrt5)/2\)，
所以
\[
 \boxed{v=c\sqrt{1-\gamma^{-2}}
 \approx0.786c}.
\]
两问不能混淆为同一个速度。

\paragraph{习题 8.14}
电子的静止能量
\(m_ec^2\simeq0.511\,\mathrm{MeV}\)；
给动能 \(K=2.8\times10^9\,\mathrm{eV}\)，
有
\[
 \gamma=1+\frac K{m_ec^2}\approx5480.46.
\]
直接计算
\(c-v=c[1-\sqrt{1-\gamma^{-2}}]\)
存在严重浮点减法消去，
应写稳定形式
\[
 \boxed{c-v=\frac{c\,\gamma^{-2}}
 {1+\sqrt{1-\gamma^{-2}}}
 \approx4.99\,\mathrm{m/s}}.
\]
能量—动量关系
\(p^2c^2=K(K+2m_ec^2)\)，
故
\[
 \boxed{p=\frac{\sqrt{K(K+2m_ec^2)}}c
 \approx1.50\times10^{-18}\,\mathrm{kg\,m/s}}
\]
或 \(pc\simeq2.80051\,\mathrm{GeV}\)。
储存环半径
\(R=L/(2\pi)=240/(2\pi)\simeq38.2\,\mathrm m\)，
横向所需向心力由动量转向率确定，
不是简单以静质量套 \(mv^2/R\)：
\[
 \boxed{F=\frac{pv}{R}\approx
 1.17\times10^{-11}\,\mathrm N}.
\]
由垂直磁场洛伦兹力
\(evB=pv/R\)，
\[
 \boxed{B=p/(eR)\approx0.245\,\mathrm T}.
\]

\paragraph{习题 8.15}
对单个宇宙线质子，
总能量按题给 \(E\simeq50\,\mathrm J\)
（由于远大于静能，将50J视为动能也可忽略差别）。
质子静能约
\(m_pc^2=1.503\times10^{-10}\,\mathrm J\)，
\(\gamma=E/(m_pc^2)\simeq3.33\times10^{11}\)。
于是
\[
 \boxed{c-v=
 \frac{c}{\gamma^2(1+\sqrt{1-\gamma^{-2}})}
 \simeq\frac{c}{2\gamma^2}
 \approx1.35\times10^{-15}\,\mathrm{m/s}}.
\]
它远低于普通双精度表示中
\(c\) 附近的绝对间隔，
必须采用上式而非直接数值相减。

\paragraph{习题 8.16}
氘核由自由质子与中子结合时，
质量亏损
\[
 \Delta m=m_p+m_n-m_D
 =(1.67265+1.67495-3.34365)\times10^{-27}
 =3.95\times10^{-30}\,\mathrm{kg}.
\]
结合能
\[
 \boxed{B_D=\Delta mc^2
 \approx2.216\,\mathrm{MeV}}.
\]
占氘核静能之比
\[
 \boxed{\frac{B_D}{m_Dc^2}\times100\%
 =\frac{\Delta m}{m_D}\times100\%
 \approx0.118\%}.
\]
氢原子的电子结合能为
\(13.58\,\mathrm{eV}\)；
氢原子静能用题给 \(m_H\)，
故
\[
 \boxed{\frac{13.58\,\mathrm{eV}}
 {m_Hc^2}\times100\%
 \approx1.45\times10^{-6}\%}.
\]
核结合所占静能比例远高于
电子与质子的原子结合比例。

\paragraph{习题 8.17}
本题用主教材给出的简化四质子聚变式，
一次反应的初质量减去终产物质量：
\[
 \begin{aligned}
 \Delta m&=4m_p-m_{\rm He}-2m_e\\
 &=[4(1.6726)-6.6425-2(0.0009)]
 \times10^{-27}\,\mathrm{kg}\\
 &=\boxed{4.61\times10^{-29}\,\mathrm{kg}}.
 \end{aligned}
\]
放出能量
\[
 \boxed{Q=\Delta mc^2\approx
 4.14\times10^{-12}\,\mathrm J
 \approx25.86\,\mathrm{MeV}}.
\]
质量转换效率
\[
 \boxed{\eta=\frac{\Delta m}{4m_p}
 \approx0.006890=0.6890\%}.
\]
每消耗 \(1\,\mathrm{kg}\) 质子的释放能
\[
 \boxed{q=\eta c^2\approx
 6.19\times10^{14}\,\mathrm{J/kg}}.
\]
按太阳辐射功率
\(P=3.9\times10^{26}\,\mathrm W\)，
\[
 \boxed{\dot m_p=P/q
 \approx6.30\times10^{11}\,\mathrm{kg/s}}.
\]
若太阳按题设有
\(M_p=1.5\times10^{30}\,\mathrm{kg}\)
可用质子，并持续此速率：
\[
 \boxed{T=M_p/\dot m_p
 \approx2.38\times10^{18}\,\mathrm s
 \approx7.55\times10^{10}\,\mathrm{yr}}.
\]
这是按源书反应式和固定太阳辐射功率
建立的粗略上限，不是太阳真实
主序寿命或氢燃烧区域的精确预测。

\paragraph{习题 8.18}
无穷远处静止的单位质量物质
落到半径 \(R\)、质量 \(M\) 的球状天体表面，
牛顿引力势能释放大小为
\(\Delta E=GMm/R\)。
取 \(m=1\,\mathrm{kg}\)，
地球质量 \(M_E\simeq5.97\times10^{24}\,\mathrm{kg}\)、
半径 \(R_E\simeq6.37\times10^6\,\mathrm m\)，
\[
 \boxed{\Delta E_E\simeq6.26\times10^7\,\mathrm J}.
\]
相对于质量静能 \(mc^2\)
的百分比
\[
 \boxed{\eta_E=GM_E/(R_Ec^2)
 \simeq6.96\times10^{-8}\%}.
\]
中子星取 \(M\simeq1.99\times10^{30}\,\mathrm{kg}\)、
半径 \(R=1.0\times10^4\,\mathrm m\)，
\[
 \boxed{\Delta E_N\simeq1.33\times10^{16}\,\mathrm J},
\qquad
 \boxed{\eta_N=GM/(Rc^2)\approx14.8\%}.
\]
与习题8.17的简化核反应效率约0.689%比较，
中子星引力释放效率高约21倍。
\textbf{适用界限：}本题遵循教材的牛顿估算，
但 \(GM/(Rc^2)\) 已不很小，
严格的中子星引力红移、入射能量与发射辐射
须用广义相对论，不能把14.8%当作精密值。

\paragraph{习题 8.19}
取两个质子各以 \(\pm0.5c\)
在共同点静止系运动，静质量均为 \(m_p\)。
\(\gamma=1/\sqrt{1-0.5^2}=2/\sqrt3\)，
各自的动量和总能量
\[
 \boxed{\boldsymbol p_{1,2}
 =\pm\frac{m_pc}{\sqrt3}\hat{\boldsymbol x}},
 \qquad
 \boxed{E_{1,2}=\frac2{\sqrt3}m_pc^2}.
\]
在质子1的静止系中，
质子2的速度用相对论合成
\[
 v_{2|1}=\frac{-0.5c-0.5c}
 {1+0.25}=-0.8c,\qquad
 \gamma_{2|1}=5/3.
\]
于是
\[
 \boxed{\boldsymbol p_{2|1}=
 -\frac43m_pc\,\hat{\boldsymbol x}},
 \quad
 \boxed{E_{2|1}=\frac53m_pc^2}.
\]
\emph{另一个}质子1在本系静止：
\(p_{1|1}=0,E_{1|1}=m_pc^2\)。
动能若需要，可从总能量中减去静能。

\paragraph{习题 8.20}
采用高能粒子惯用单位，
初态是总能量为 \(E_e\) 的运动电子，
与静止质子 \(m_pc^2\) 相撞。
实验室系初四动量求质心不变量
\[
 s c^4=m_e^2c^4+m_p^2c^4
 +2m_pc^2E_e.
\]
题给电子\emph{动能}22GeV，
因此 \(E_e=22\,\mathrm{GeV}+m_ec^2\)。
用
\(m_pc^2\simeq0.938272\,\mathrm{GeV}\)，
\(m_ec^2\simeq0.000511\,\mathrm{GeV}\)，
得到
\[
 \boxed{\sqrt{s}\,c^2
 \approx6.4935\,\mathrm{GeV}},
\]
即质心系全部可利用总能量。
若「资用能」专指\emph{扣除初态两粒子静能}
之后可转成新粒子静能/动能的部分，
则
\[
 \boxed{E_{\rm available,kin}
 =\sqrt{s}\,c^2-(m_p+m_e)c^2
 \approx5.5547\,\mathrm{GeV}}.
\]
两种定义在教材不同术语体系中可能有区别，
须与原第8章「资用能」的定义逐式校对，
不可只报一个数又不说明是否扣除静能。

\paragraph{习题 8.21}
相向对撞的电子、正电子
两者能量各约 \(2.8\,\mathrm{GeV}\)，
初始动量相反，质心系即实验室系。
故总质心能量
\[
 \boxed{E_{\rm cm}\approx5.6\,\mathrm{GeV}}.
\]
若「资用能」扣除两个电子静能，
则约
\(5.6-2(0.000511)=5.5990\,\mathrm{GeV}\)。
对于同样电子射向\emph{静止电子}，
要达到同一质心总能量，
\[
 E_{\rm cm}^2=2m_e^2c^4+2m_ec^2 E_{\rm beam},
\]
故实验室系\emph{入射电子总能量}
\[
 \boxed{E_{\rm beam}=
 \frac{E_{\rm cm}^2-2m_e^2c^4}
 {2m_ec^2}
 \approx3.07\times10^4\,\mathrm{GeV}}.
\]
入射电子动能仅少
\(0.000511\,\mathrm{GeV}\)，
也约为 \(3.07\times10^4\,\mathrm{GeV}\)。
与对撞机单束2.8GeV相比，
固定靶所需束能约高
\(\boxed{1.10\times10^4}\) 倍，
说明同等束能时相向对撞更有效。

\paragraph{习题 8.22}
沿 \(x\) 的洛伦兹四动量变换：
\[
 E'=\gamma(E-up_x),\quad
 p_x'=\gamma(p_x-uE/c^2),\quad
 p_y'=p_y,\quad p_z'=p_z.
\]
直接展开：
\[
 \begin{aligned}
 E'^2-c^2p_x'^2
 &=\gamma^2[(E-up_x)^2
 -c^2(p_x-uE/c^2)^2]\\
 &=\gamma^2(1-u^2/c^2)
 (E^2-c^2p_x^2)\\
 &=E^2-c^2p_x^2.
 \end{aligned}
\]
再减去未变的
\(c^2p_y'^2+c^2p_z'^2\)，得到
\[
 \boxed{E'^2-c^2|\boldsymbol p'|^2
 =E^2-c^2|\boldsymbol p|^2
 =m_0^2c^4}.
\]
不变量不依赖所选惯性参考系。

\paragraph{习题 8.23}
静止质子的静能
\(E_0=m_pc^2\)，
\[
 \boxed{E_0\simeq0.938\,\mathrm{GeV}}.
\]
代入牛顿引力形式：
两个静止质子间
\(F_g=G(E_0/c^2)^2/r^2
 =GE_0^2/(c^4r^2)\)。
题目\emph{两质子均具有相同高能 \(E\)}
且粗略把 \(E/c^2\) 当各自的引力等效质量，
则
\[
 \frac{F_g(E)}{F_g(E_0)}
 =\left(\frac{E}{E_0}\right)^2.
\]
欲令原本比引力大约
\(10^{36}\) 倍的电力被引力追平，
需
\[
 \left(\frac{E}{E_0}\right)^2\simeq10^{36}
 \Longrightarrow
 \boxed{E\simeq10^{18}E_0
 \approx9.38\times10^{17}\,\mathrm{GeV}}.
\]
比较教材给出的粒子加速器量级
\(10^3\,\mathrm{GeV}\)：
\[
 \boxed{\frac{E}{10^3\,\mathrm{GeV}}
 \approx9.38\times10^{14}}.
\]
\textbf{模型界限：}若只把\emph{一个}
质子的能量提高、另一质子保持静止，
比例会只与 \(E/E_0\) 成正比，
结果不同；主教材题目默认讨论
「两个质子」共同的典型能标。
把动能与重力质量简单代入牛顿
\(1/r^2\) 定律并非严格量子引力理论，
这里只用于揭示约 \(10^{18}\) GeV 的极高能标。
